\subsection{Graph Reduction: Definitions}
\label{sec:prelim:reduction}
A reduction operator $R$ maps a graph $G$ to a smaller graph $R(G)$ intended to stand in
for it. The taxonomy used throughout follows the survey of \citet{hashemi2024survey},
which divides the field into sparsification, coarsening and condensation, with one
substitution: partitioning replaces condensation as the third family. Partitioning is
what a distributed training setup forces on a graph that does not fit one device
(\ref{sec:relwork:scaling:distributed}), whereas condensation is excluded from this study
outright, for the reasons given in~\ref{sec:relwork:reduction:summarization}. The
substitution is deliberate, not an error of transcription.

\subsubsection{Partitioning}
\label{sec:prelim:reduction:partitioning}
Partitioning splits $\mathcal{V}$ into $k$ disjoint blocks and deletes every edge whose
endpoints fall in different blocks. It is the one family that keeps every node: only
spanning edges are lost, so the node count of $R(G)$ always equals that of $G$ and any
compression it achieves is compression of the edge set alone. What varies between
partitioners is the objective by which blocks are chosen
(\ref{sec:method:reduction:partitioning}).

\subsubsection{Sparsification}
\label{sec:prelim:reduction:sparsification}
Sparsification removes edges, or nodes together with their incident edges, while
preserving some structural property of the original, such as pairwise distances up to a
stretch factor. The two variants are not interchangeable for measurement: an edge-removing
method leaves $|\mathcal{V}|$ untouched while a node-removing method shrinks it, and the
distinction is what forces the two-ratio convention
of~\ref{sec:prelim:reduction:measuring}.

\subsubsection{Summarization and Coarsening}
\label{sec:prelim:reduction:summarization}
Summarization merges nodes into \emph{super-nodes}. A coarsening is specified by a merge
map, a surjection $\mu \colon \mathcal{V} \rightarrow \mathcal{V}'$ onto the cluster set,
and the resulting quotient graph connects two clusters whenever any of their members were
connected. The number of member edges a super-edge stands for is its multiplicity,
\begin{equation}
    w_{u'v'} = \bigl| \{\, (u, v) \in \mathcal{E} :
        \mu(u) = u',\ \mu(v) = v' \,\} \bigr|,
    \label{eq:prelim:multiplicity}
\end{equation}
and when super-edges differing in an attribute are kept parallel rather than collapsed,
the quotient is a multigraph. Because $\mu$ is a function, its preimages partition
$\mathcal{V}$: a coarsening \emph{is} a partition of the node set. A tolerance band such
as ``within one level'' therefore cannot define one, since the relation it induces is not
transitive, a fact the domain-specific candidate of~\ref{sec:method:summary:domain} has
to work around.

\subsubsection{Measuring Reduction: Compression Ratio}
\label{sec:prelim:reduction:measuring}
Reduction is reported as retention, kept over original, separately per node and per edge:
\begin{equation}
    \eta_{\mathcal{V}}(G) =
        \frac{\left| \mathcal{V}(R(G)) \right|}{\left| \mathcal{V}(G) \right|},
    \qquad
    \eta_{\mathcal{E}}(G) =
        \frac{\left| \mathcal{E}(R(G)) \right|}{\left| \mathcal{E}(G) \right|},
    \label{eq:prelim:retention}
\end{equation}
with reduction $1 - \eta$. One number cannot carry both: partitioning and edge-removing
sparsification hold $\eta_{\mathcal{V}} = 1$ while lowering $\eta_{\mathcal{E}}$, and
node-removing methods move the two together. Both ratios are therefore reported
everywhere, and never collapsed (\ref{sec:method:experiment:metrics:reduction}). Methods
further differ in whether compression is a knob or an outcome: some accept a target
ratio, while for others $\eta$ is a fixed property of the input graph. Comparing methods
at matched compression is consequently a design problem rather than a reporting
convention, and~\ref{sec:method:summary:random} is its solution.
